\documentclass[17pt,a4paper,landscape]{extarticle}
\usepackage[margin=2.35cm,top=0.12cm,left=1.05cm]{geometry}
\usepackage{amsmath}
\usepackage{amssymb}
\usepackage{tasks}
\usepackage{tikz}
\usepackage{xcolor}
\usepackage[sfdefault,lf]{FiraSans}
\renewcommand{\familydefault}{\sfdefault}
\setlength{\parindent}{0pt}
\pagestyle{empty}
\settasks{label=\textbf{\arabic*.}, label-width=2.2em, item-indent=2.95em, column-sep=2.1cm, after-item-skip=4.7em}
\renewcommand{\arraystretch}{1.15}
\everymath{\displaystyle}

\begin{document}
\sffamily
\boldmath

\boldmath

\boldmath

\boldmath

\boldmath

\boldmath

\boldmath

{\LARGE \textbf{Proficient}}
\begin{tasks}(3)
\task After a 15\% increase (GST), a product costs \$115. What was the price before GST?
\task After a 20\% increase, a salary is \$4800. What was the original salary?
\task After a 5\% increase, a house's value is \$315,000. What was the original value?
\task After a 12.5\% increase, a ticket costs \$45. What was the original cost?
\task After a 30\% increase, a population is 6500. What was the original population?
\task After a 8\% increase, a bill is \$162. What was the original bill?
\task A price increased by 20\% then decreased by 10\% to become \$216. What was the original price?
\task A number increased by 15\% then decreased by 5\% to become 218.5. What was the original number?
\task After a 25\% discount followed by a 10\% discount, an item costs \$135. What was the original price?
\task A salary increased by 10\% then increased by another 10\% to become \$6050. What was the original salary?
\task After a 20\% discount then a 5\% discount, a TV costs \$380. What was the original price?
\task A stock price decreased by 30\% then increased by 40\% to become \$98. What was the original price?
\end{tasks}

\end{document}